\scanpage{205}
\mspage{14}
% archive sheet number: 204

décomposition est induite par une décomposition de $X$. La question étant locale sur $S$, on peut supposer $\exists\,M\in\mathrm{Ob}\,C_S$, $\varphi(M)\longrightarrow e$ un épimorphisme.

Considérons $r(X\times\varphi(M))=\alpha(X)(M)\inserted{=Z(M)}\longrightarrow r\varphi(M)=s(M)$ il se coupe en $Z_1(M)\amalg Z_2(M)$, donc \deleted{par} R2 $X\times\varphi(M)$ se coupe en $X_1\amalg X_2$. Je dis que

\diagram{\begin{tikzcd}[ampersand replacement=\&,column sep=small,row sep=small]
T\arrow[d,shift left=.5ex]\arrow[d,shift right=.5ex]\&X\times T\arrow[d,shift left=.5ex]\arrow[d,shift right=.5ex]\&\\
\varphi(M)\arrow[d]\&X\times\varphi(M)\arrow[l,dash]\arrow[d]\&X_1\arrow[l,hook]\\
e\&X\arrow[l,dash]\&
\end{tikzcd}}

cette décomposition provient d'une décomposition de $X$, i.e. que quels que soient \deleted{\unsure{du|les}} l'objet \deleted{\unsure{Z|T}}\inserted{$T$} de $R_S^{C_0}$ et deux morphismes $T\rightrightarrows\varphi(M)$, les deux images inverses de $X_1$ par $T\rightrightarrows\varphi(M)$ sont les \unsure{mêmes}, ou

% slanted
\margin{\unsure{encore dont|encore deux|encore dans|on a} les $\leftleftarrows$ \unsure{sont} \deleted{\unsure{égaux}} $r(X_1)\rightrightarrows r(X\times T)$}

Or, \inserted{quitte à ajouter à $T$ un objet $\varphi(M)$ (qui envoie deux fois dans $\varphi(M)$ par id) on peut supposer} \deleted{quitte à localiser $T$ lui est subordonné à} \deleted{$\varphi(M)$, ce qui implique} $T\longrightarrow e_S$ épim., donc $T\times\varphi(M)\longrightarrow\varphi(M)$ un épim. Comme $T\deleted{\times\varphi(M)}$ est subordonné à $\varphi(M)$,

$T\times\varphi(M)$ est constant par morceaux sur $\varphi(M)$, et étant épimorphique, \unsure{cela|il} est \uline{recouvert} (pour sa topologie canonique, p. ex.) par les images des sections $w$ de \deleted{\unsure{$\varphi$}} $T\times\varphi(M)$ sur $\varphi(M)$, Ceci posé, on peut remplacer $T\rightrightarrows\varphi(M)$ par les composés $T\times\varphi(M)\xrightarrow{\mathrm{pr}_1}T\rightrightarrows\varphi(M)$ (car $\mathrm{pr}_1$ est épim., $\varphi(M)\longrightarrow e_S$ l'étant !), puis \unsure{ces derniers|ce dernier|en dernier} par les composés $\varphi(M)\xrightarrow{w}T\times\varphi(M)$. Bref, on est ramené au cas où $T=\varphi(M)$. Mais $\varphi$ étant pl. fidèle

